\documentclass[11pt,a4paper]{article}
\usepackage[utf8]{inputenc}
\usepackage{amsmath,amssymb,amsthm}
\usepackage{geometry}
\geometry{margin=2.5cm}

\newtheorem{theorem}{Theorem}[section]
\newtheorem{lemma}[theorem]{Lemma}
\theoremstyle{definition}
\newtheorem{definition}[theorem]{Definition}

\newcommand{\N}{\mathbb{N}}
\newcommand{\R}{\mathbb{R}}
\newcommand{\C}{\mathbb{C}}
\newcommand{\Z}{\mathcal{Z}}          % non‑trivial zeros of ζ
\newcommand{\LOp}{\mathcal{L}}
\newcommand{\TOp}{T}
\newcommand{\supp}{\operatorname{supp}}
\newcommand{\spec}{\operatorname{spec}}
\DeclareMathOperator{\Ree}{Re}
\DeclareMathOperator{\Imm}{Im}

\title{Summary and Critical Gap in the GDST Proof of the Riemann Hypothesis}
\author{}
\date{}

\begin{document}
\maketitle

\begin{abstract}
We review the core structure of the claimed proof of the Riemann Hypothesis via the Greedy Dirichlet Sub‑Sum Transform (GDST) and identify the point at which the argument fails.  The proposed derivation of a central identity linking the regularised Stieltjes transform of the spectral measure of \(T_\sigma\) to a sum over the non‑trivial zeros of \(\zeta\) is not mathematically justified, because it ignores the conditional convergence of the zero‑sum and omits essential regularising terms.  Consequently the whole proof collapses.
\end{abstract}

\section{Structure of the GDST proof}
The proof is organised in five parts.
\begin{enumerate}
  \item A greedy expansion of real numbers into harmonic fractions \(\alpha_n = \frac{\pi}{n+2}\) yields a family of binary selection indicators \(\delta_n(\theta)\in\{0,1\}\).
  \item The correlation kernel \(K(n,m)=\int_0^\pi[\delta_n(\theta)-\theta/\pi][\delta_m(\theta)-\theta/\pi]\,d\theta\) is shown to be positive semi‑definite.
  \item For \(\Re s>0\) the GDST Dirichlet series 
    \[
    E_{\GDST}(\theta,s)=\sum_{n=0}^\infty \delta_n(\theta)(n+2)^{-s}
    \]
    converges absolutely and is linked to \(\zeta(s)-1\).
  \item A transfer operator \(\LOp_s\) on the Hardy space \(H^2(\mathbb{D})\) is introduced, and its Fredholm determinant is expressed as
    \[
    {\det}_{(2)}(I-\LOp_s)=\frac{\zeta(s)}{\zeta(2s)}e^{P(s)},
    \]
    where \(P(s)\) is entire.
  \item For a real parameter \(\sigma>0\) a self‑adjoint positive trace‑class operator \(T_\sigma\) is defined with purely discrete spectrum \(\lambda_i(\sigma)\ge0\).  Its spectral measure is
    \[
    \mu_\sigma = \sum_i \delta_{\lambda_i(\sigma)}.
    \]
    The regularised Stieltjes transform
    \[
    S_{\mathrm{reg}}(z)=\sum_i\Bigl(\frac1{z-\lambda_i(\sigma)}-\frac1z\Bigr)
    \]
    is introduced, and its moments \(M_k(\sigma)=\sum_i\lambda_i(\sigma)^k\) are expressed, via trace identities and the determinant formula of Part~IV, as
    \[
    M_k(\sigma)=\sum_{\rho\in\Z}\frac1{(\rho-\sigma)^k}+H_k(\sigma),\qquad |H_k(\sigma)|\le C A^k,
    \]
    where \(\Z\) denotes the set of non‑trivial zeros of \(\zeta\).
\end{enumerate}

\section{The attempted pole matching}
From the moment expansion one writes
\[
S_{\mathrm{reg}}(z)=\sum_{k=1}^\infty\frac{M_k(\sigma)}{z^{k+1}},
\qquad |z|>R:=\sup_i\lambda_i(\sigma).
\]
Substituting the trace identity and interchanging summations (justified by the geometric nature of the error series and a pairing of symmetric zeros \(\rho\) and \(1-\rho\)) the authors claim that
\begin{equation}\label{eq:claim}
S_{\mathrm{reg}}(z) = \sum_{\rho\in\Z}\frac1{z-(\rho-\sigma)} + \Phi(z),
\end{equation}
where \(\Phi\) is an entire function.

If \eqref{eq:claim} were true, the meromorphic function \(S_{\mathrm{reg}}\) would have simple poles exactly at \(\rho-\sigma\).  Because the left‑hand side has poles precisely at the real eigenvalues \(\lambda_i(\sigma)\), one would deduce
\[
\{\lambda_i(\sigma)\}=\{\rho-\sigma:\rho\in\Z\},
\]
so every \(\rho-\sigma\) is real.  Choosing \(\sigma=\frac12\) would then force \(\Im\rho=0\) for all non‑trivial zeros, and by combining this with the functional equation and the properties of \(\zeta\) one would conclude \(\Re\rho=\frac12\).

\section{The fatal gap}
The identity \eqref{eq:claim} is **not** derived correctly.  The series \(\sum_{\rho}1/(\rho-\sigma)^k\) is only conditionally convergent for \(k=1\), and the sum over zeros of \(1/(z-(\rho-\sigma))\) does **not** converge absolutely, nor does it define a meromorphic function with simple poles at the shifted zeros without a suitable regularising term (such as the additional \(1/\rho\) present in the standard product formula for \(\zeta\)).

The algebraic manipulation that attempts to produce \eqref{eq:claim} essentially writes
\[
\sum_{k=1}^\infty\frac{1}{z^{k+1}}\sum_{\rho}\frac1{(\rho-\sigma)^k}
= \frac12\sum_{\rho}\sum_{k=1}^\infty\frac{P_k(\rho)}{z^{k+1}},
\]
with \(P_k(\rho)=\frac1{(\rho-\sigma)^k}+\frac1{(1-\rho-\sigma)^k}\).  After summing the inner geometric series one obtains
\[
\sum_{k=1}^\infty\frac{P_k(\rho)}{z^{k+1}}
= \frac1{z-(\rho-\sigma)}+\frac1{z-(1-\rho-\sigma)}-\frac2z .
\]
Summing over all \(\rho\) and using the symmetry \(\rho\leftrightarrow1-\rho\) yields
\[
\sum_{\rho\in\Z}\Bigl(\frac1{z-(\rho-\sigma)}-\frac1z\Bigr).
\]
The term \(\sum_\rho\frac1z\) diverges, because the sum \(\sum_\rho 1\) is infinite.  The authors claim that this divergent constant ``cancels against a regularising entire function''.  No such entire function is constructed; the only entire function in the derivation, \(\Phi(z)\), arises from the error terms \(H_k(\sigma)\) and has a power series with geometrically decaying coefficients — it absolutely cannot cancel an infinite constant.  The cancellation is therefore an algebraic oversight, not a legitimate analytic step.

In the standard theory of the zeta function, the corresponding regularisation is provided by adding \(\frac1\rho\) inside the sum, producing the convergent series
\[
\sum_{\rho}\Bigl(\frac1{s-\rho}+\frac1\rho\Bigr)
= \frac{\zeta'(s)}{\zeta(s)} + (\text{elementary terms}).
\]
The GDST manuscript does **not** incorporate this regularisation, and no alternative mechanism is supplied.  Hence \eqref{eq:claim} is unsubstantiated, and the entire logical chain that leads to the Riemann Hypothesis collapses.

\section{What remains}
The remaining gap is fundamental: to prove the Riemann Hypothesis via this operator‑theoretic route one would need to establish a rigorous link between the spectral zeta function of \(T_\sigma\) and the zero set of \(\zeta\), a link that goes far beyond matching formal power series.  The techniques presented in the manuscript do not achieve this, and the additional documents examined during this discussion do not contain the necessary breakthrough.  The claimed proof of RH is therefore incorrect.
\end{document}